\section{Subobjects, negation, and two laws of logic}\label{sec:negation}

\subsection{Sub-presheaves}

\begin{definition}[Sub-presheaf]
A \emph{sub-presheaf} $A$ of a presheaf $F$ is a choice of subsets $A(U)\subseteq F(U)$, closed under restriction: if $x\in A(U)$ and $f : V\to U$, then $\res xf\in A(V)$.
\end{definition}

For a right $M$-set $X$, a sub-presheaf is a subset closed under the action. For an automaton, it is a set of states closed under transitions: once you are in it, you cannot get out. For a graph, it is a \emph{subgraph}: a set of vertices and a set of edges such that the endpoints of every chosen edge are chosen. For a simplicial set, it is a simplicial subset. Sub-presheaves are the subobjects of $F$ in the category $\PSh(C)$, and we use both words.

The sub-presheaves of $F$ are ordered by inclusion, and they form a lattice: the union $A\cup B$ and the intersection $A\cap B$, computed at each object, are again sub-presheaves. The smallest one is $\emptyset$, the largest is $F$ itself. Two sub-presheaves are \emph{disjoint} if $A\cap B = \emptyset$.

\subsection{Complement and pseudocomplement}

In a set, every subset has a complement. For sub-presheaves this fails: the set-theoretic complement of $A$ is usually not closed under restriction.

\begin{example}\label{ex:edge}
Let $F = \y E$ be the graph with one edge $e$ from $u$ to $v$, and $A = \{u\}$ the subgraph consisting of the vertex $u$ alone. Its set-theoretic complement $\{v, e\}$ is not a subgraph, since the edge $e$ has its source $u$ outside it.
\end{example}

\begin{figure}[ht]
\centering
\begin{tikzpicture}[v/.style={circle,fill,inner sep=1.6pt},>=Stealth]
\node[v,label=below:$u$] (u) at (0,0) {}; \node[v,label=below:$v$] (v) at (2,0) {};
\draw[->,thick] (u) -- node[above]{$e$} (v);
\node at (1,-1) {$F = \y E$};
\begin{scope}[xshift=3.6cm]
\node[v,red,label=below:$u$] (u) at (0,0) {}; \node[v,black!25] (v) at (2,0) {};
\draw[->,black!25] (u) -- (v);
\node at (1,-1) {$A = \{u\}$};
\end{scope}
\begin{scope}[xshift=7.2cm]
\node[v,black!25] (u) at (0,0) {}; \node[v,blue,label=below:$v$] (v) at (2,0) {};
\draw[->,black!25] (u) -- (v);
\node at (1,-1) {$\neg A = \{v\}$};
\end{scope}
\begin{scope}[xshift=10.8cm]
\node[v,red,label=below:$u$] (u) at (0,0) {}; \node[v,black!25] (v) at (2,0) {};
\draw[->,black!25] (u) -- (v);
\node at (1,-1) {$\neg\neg A = \{u\}$};
\end{scope}
\end{tikzpicture}
\caption{In the graph with one edge, $\neg A\cup\neg\neg A = \{u,v\}$ misses the edge.}\label{fig:edge}
\end{figure}

What survives is the largest sub-presheaf disjoint from $A$.

\begin{definition}[Pseudocomplement]
The \emph{pseudocomplement} $\neg A$ of a sub-presheaf $A\subseteq F$ is the largest sub-presheaf of $F$ disjoint from $A$.
\end{definition}

\begin{lemma}\label{lem:neg}
$\neg A$ exists, and an element $x\in F(U)$ lies in $\neg A$ iff no restriction of $x$ lies in $A$: $\res xf\notin A$ for every arrow $f$ into $U$. Consequently $x\in\neg\neg A$ iff every restriction of $x$ has a further restriction in $A$: for every $f$ into $U$ there is $g$ with $\res{(\res xf)}{g}\in A$.
\end{lemma}

\begin{proof}
Let $N(U)$ be the set of $x\in F(U)$ none of whose restrictions lies in $A$. It is closed under restriction (a restriction of $\res xf$ is a restriction of $x$), and disjoint from $A$ (take $f = \id$). If $B$ is a sub-presheaf disjoint from $A$ and $x\in B(U)$, every restriction of $x$ lies in $B$, hence not in $A$, so $x\in N(U)$. So $N$ is the largest sub-presheaf disjoint from $A$. The description of $\neg\neg A$ is the same statement applied to $\neg A$: $x\in\neg\neg A$ iff no restriction $\res xf$ lies in $\neg A$, iff each $\res xf$ has some restriction in $A$.
\end{proof}

\begin{example}[Example~\ref{ex:edge} continued]
In $F = \y E$ with $A = \{u\}$: the edge $e$ restricts to $u\in A$, so $e\notin\neg A$; the vertex $v$ has no restriction other than itself, so $v\in\neg A$; hence $\neg A = \{v\}$. By symmetry $\neg\neg A = \neg\{v\} = \{u\}$. So $\neg A\cup\neg\neg A = \{u,v\}$ is \emph{not} the whole graph: the edge is in neither (Figure~\ref{fig:edge}).
\end{example}

\begin{example}[Automata]\label{ex:automataneg}
Let $X$ be an automaton and $A$ a closed set of states (a \q{trap}). By Lemma~\ref{lem:neg}, $x\in\neg A$ iff no word leads from $x$ into $A$: \q{$x$ can never fall into the trap}. And $x\in\neg\neg A$ iff from every state reachable from $x$, some word leads into $A$: \q{from $x$, the trap remains reachable whatever you do}.
\end{example}

\begin{example}[$G$-sets]
If $G$ is a group and $A$ a sub-$G$-set of $X$, then $X\smallsetminus A$ is a sub-$G$-set: if $x\notin A$ and $x\cdot g\in A$ then $x = (x\cdot g)\cdot g^{-1}\in A$. So $\neg A = X\smallsetminus A$ and $\neg\neg A = A$: in $G$-sets, negation behaves as for sets.
\end{example}

\subsection{Boolean and De Morgan}

In classical logic, every proposition $p$ satisfies the \emph{law of excluded middle} $p\vee\neg p$; in set theory, $A\cup(X\smallsetminus A) = X$. In a category of presheaves, sub-presheaves play the role of propositions, $\cup$ of \q{or}, and $\neg$ of \q{not}. The two laws below are the ones our bridges are about.

\begin{definition}[Boolean, De Morgan]
$\PSh(C)$ is \emph{Boolean} if every sub-presheaf $A$ of every presheaf $F$ has a complement: a sub-presheaf $B$ with $A\cap B=\emptyset$ and $A\cup B = F$. Equivalently, $A\cup\neg A = F$ for all $A$ (\emph{excluded middle}).
$\PSh(C)$ is \emph{De Morgan} if $\neg A\cup\neg\neg A = F$ for every sub-presheaf $A$ of every presheaf $F$ (\emph{weak excluded middle}).
\end{definition}

(If $B$ is a complement of $A$, then $B\subseteq\neg A$ since $B$ is disjoint from $A$, so $A\cup\neg A\supseteq A\cup B = F$; conversely $\neg A$ is then a complement. So the two forms of \q{Boolean} agree.)

The name \q{De Morgan} comes from the fact that, for the intuitionistic logic of a topos, $\neg p\vee\neg\neg p$ is equivalent to the De Morgan law $\neg(p\wedge q) = \neg p\vee\neg q$ \cite{johnstone-demorgan}. Every Boolean topos is De Morgan, since $\neg\neg A = A$ when $A$ has a complement. By Lemma~\ref{lem:neg}:

\begin{lemma}\label{lem:dmpointwise}
$\PSh(C)$ is De Morgan iff, for every presheaf $F$, every sub-presheaf $A$ and every element $x$ of $F$: either no restriction of $x$ lies in $A$, or every restriction of $x$ has a further restriction in $A$.
\end{lemma}

So, for automata: the topos of $\Sigma$-automata is De Morgan iff, in every automaton and for every trap, every state \emph{either can never fall into the trap, or can never escape the possibility of falling into it}. Whether this holds depends on $\Sigma$, as we shall see in \S\ref{sec:monoids}.

\begin{example}[Simplicial sets]\label{ex:Delta1}
Let $F = \Delta[1] = \y[1]$, the standard $1$-simplex: two vertices $0,1$ and an edge between them (plus degenerate simplices). Take $A$ = the vertex $0$ (with its degeneracies). Exactly as for the one-edge graph, $\neg A$ is the vertex $1$, $\neg\neg A$ is the vertex $0$, and $\neg A\cup\neg\neg A = \partial\Delta[1]$, the two endpoints without the edge. So simplicial sets do not form a De Morgan topos.
\end{example}

\begin{example}[The topos of trees]\label{ex:trees}
Presheaves on $\omega$: sequences $X_0\xleftarrow{r}X_1\xleftarrow{r}X_2\leftarrow\cdots$. An element $x\in X_n$ has restrictions $r^k(x)\in X_{n-k}$. Let $A$ be a sub-presheaf and $x\in X_n$. If some restriction $r^k(x)$ lies in $A$, then every restriction $r^j(x)$ has a further restriction in $A$: namely $r^{\max(j,k)}(x)$, which is a restriction of $r^k(x)$, hence in $A$. So the topos of trees is De Morgan. It is not Boolean: in $X = \y 1$ (that is $X_0 = X_1 = \{\ast\}$, $X_n = \emptyset$ for $n\ge 2$), the sub-presheaf $A$ with $A_0 = \{\ast\}$, $A_1 = \emptyset$ has $\neg A = \emptyset$.
\end{example}

\begin{exercise}\label{exo:heyting}
Show that, for sub-presheaves, $A\subseteq\neg\neg A$, $\neg\neg\neg A = \neg A$, and $\neg(A\cup B) = \neg A\cap\neg B$. (These hold in every topos; the lattice of sub-presheaves is a \emph{Heyting algebra}.)
\end{exercise}

\begin{exercise}\label{exo:dmlaw}
Show that $\PSh(C)$ is De Morgan iff $\neg(A\cap B) = \neg A\cup\neg B$ for all sub-presheaves $A$, $B$ of every presheaf.
\end{exercise}
